\section{Protocol and Setup}
\label{sec:setup}

\subsection{Data and splits}

We use the four C-MAPSS subsets \cite{saxena2008}: FD001 (one condition, one
fault mode), FD002 (six conditions, one fault mode), FD003 (one condition, two
fault modes), and FD004 (six conditions, two fault modes). The RUL target is
capped at 125 cycles. Engines are split into training and validation engines;
validation engines are never used for fitting, early stopping, or model
selection. The official test engines are sealed until the frozen configuration
is fixed; we do not use test labels for any choice in this paper.

For validation we evaluate five cut fractions of each engine's life:
0.40, 0.55, 0.70, 0.85, and 0.95. Each cut produces one prediction per engine.
The official test provides one window per engine near end-of-life, so the 0.95
cut is the closest validation analogue. This distinction is central to the
analysis in Section~\ref{sec:cutwise}.

\subsection{Backbone and baselines}

The backbone is a frozen DSPT-style model (B0) trained under a fixed protocol:
two layers, model width 64, dropout 0.2, FP32, 40 epochs, patience 12, batch
size 512, learning rate $10^{-3}$, and weight decay $10^{-4}$. We do not retrain
the backbone for any experiment in this paper. The single-window prediction from
this backbone is denoted K1. The median of age-aligned windows is denoted Median.

\subsection{Estimators}

Two estimators are reported and never mixed:
\begin{itemize}
\item \emph{Engine-mean RMSE}: compute RMSE per engine, then average over
engines. This gives every engine equal weight.
\item \emph{Pooled RMSE}: pool all windows and compute one RMSE. This weights
engines by their number of windows.
\end{itemize}
The identity $\mathrm{pooled}^2-\mathrm{engine\_mean}^2=\mathrm{Var}_e(\mathrm{RMSE}_e)$
holds when every engine has the same number of windows; we use it as a
consistency check. The official test has one window per engine, so the two
estimators coincide there.

For paired comparisons we resample engines with replacement (10{,}000 draws) and
report the mean difference, its 95\% confidence interval, and a one-sided
$p$-value. When multiple seeds are available, predictions are averaged within an
engine before resampling, so the unit of replication is the engine, not the seed.

\subsection{Reproducibility note}

The Phase~2 experiments in the broader project used a worktree code lock rather
than the original commit; the lock file and the current worktree hashes are
retained in the repository. The LTAG experiments in this paper use the frozen
B0 checkpoint and the feature table described in Section~\ref{sec:method}; no
test-time tuning is performed.
